\subsection{完全域的扩张}

为便于引用，我们总结完全域扩张的主要性质：

定理 3. --- 设 K 为完全域。

\begin{enumerate}
\def\labelenumi{\alph{enumi})}
\item
  \(\mathbf{K}\) 的每个代数扩张都是完全域。
\item
  \(\mathbf{K}\) 的每个扩张都是可分的。
\item
  一个 K-代数可分当且仅当它是既约的。
\item
  设 \(\mathbf{A}\) 和 \(\mathbf{B}\) 是两个既约 K-代数。则 \(\mathbf{A} \otimes_{\mathbf{K}} \mathbf{B}\) 是既约的。
\item
  如果 \(\mathbf{E}\) 和 \(\mathbf{F}\) 是 \(\mathbf{K}\) 的两个扩张，则环 \(\mathbf{E} \otimes_{\mathbf{K}} \mathbf{F}\) 是既约的。
\end{enumerate}

断言 \(a\) ) 正是命题 11 的推论 1（第 V 卷，第 43 页）。

当 K 的特征为 0 时，断言 b) 由定理 1（V，第 122 页）得出；当 K 的特征为 \(p \neq 0\) 时，则由推论 1（V，第 123 页）得出。

让我们证明 \(c\) )。\(\mathbf{K}\) 的特征为 0 的情形由定理 1（V，第 122 页）得出。因此只需证明：如果 \(\mathbf{K}\) 是特征为 \(p \neq 0\) 的完全域，且 \(\mathbf{A}\) 是既约 K-代数，则 \(\mathbf{A}\) 在 \(\mathbf{K}\) 上可分；但这由定理 2 的条件 \(a\) ) 与 \(b\) ) 的等价性得出（V，第 122 页；取 \(\mathbf{K}' = \mathbf{K}\) 于 \(b\) ) 中）。

最后， \(d)\) 由 \(c)\) 与命题 5（V，第 120 页）得出，而 \(e)\) 是 \(d)\) 的特殊情形。

